\section{Related Work}
\label{sec:related}

\subsection{KV Cache Quantization}

\paragraph{Uniform quantization.}
KIVI~\citep{kivi2024} applies 2-bit quantization to both Keys and Values, achieving $\sim$8$\times$ compression but with 1--3\% perplexity degradation.
KVQuant~\citep{kvquant2024} uses per-channel quantization with outlier handling, improving quality at 4-bit.
Both methods use uniform bit allocation across Keys and Values, missing the opportunity for asymmetric allocation.

\paragraph{Rotation-based quantization.}
QuIP\#~\citep{quip2024} and TurboQuant apply Hadamard rotation before quantization to reduce outlier sensitivity.
The rotation makes entries approximately i.i.d.\ Gaussian, enabling near-optimal Lloyd-Max scalar quantization.
QJL~\citep{qjl2024} uses Johnson-Lindenstrauss projections for Key compression specifically, providing theoretical guarantees via the JL lemma.
Our work builds on the TurboQuant rotation framework but extends it with asymmetric K/V allocation and integration with eviction.

\paragraph{Asymmetric quantization.}
KVTuner~\citep{kvtuner2025} (ICML 2025) studies the sensitivity difference between Keys and Values, concluding that Keys are $\sim$2$\times$ more sensitive due to softmax nonlinearity.
Our experiments reveal a more nuanced picture: while Keys have higher \emph{sensitivity} (per-bit impact on attention logits), Values have a sharper \emph{quality cliff} between 2-bit and 4-bit due to the limited codebook size at very low bit-widths.
The practical implication is that K6V4 (Keys at 6-bit, Values at 4-bit) outperforms uniform 5-bit allocation.

\subsection{KV Cache Eviction}

\paragraph{Attention-based eviction.}
H\textsubscript{2}O~\citep{h2o2023} identifies ``heavy hitter'' tokens that accumulate high attention scores and retains only these, achieving $\sim$5$\times$ compression.
ScissorHands~\citep{scissorhands2024} uses a similar importance-based criterion with a pivot token mechanism.
Both methods achieve moderate compression but lack theoretical error bounds connecting eviction to quantization noise.

\paragraph{Streaming and windowed approaches.}
StreamingLLM~\citep{streamingllm2024} observes that the first few tokens (``attention sinks'') accumulate disproportionate attention regardless of content, and proposes retaining these plus a sliding window of recent tokens.
Our eviction strategy builds on this observation but derives the eviction threshold from quantization error analysis rather than using a fixed window size.

\paragraph{Token merging.}
Token merging~\citep{tome2023} in vision transformers combines similar tokens to reduce sequence length.
Attention Matching (MIT) applies a similar idea to KV caches, replacing $n$ KV entries with $m \ll n$ virtual entries optimized to preserve attention patterns.
We discuss token merging as a future direction (Section~\ref{sec:phase3}) that could replace hard eviction with soft merging.

\subsection{Combined Approaches}

GEAR~\citep{gear2024} combines low-rank approximation with sparse outlier correction, achieving 4--8$\times$ compression.
However, the low-rank assumption does not hold uniformly across layers, limiting compression ratio.
Our approach differs by combining quantization (which has strong theoretical guarantees via rate-distortion theory) with eviction (which exploits the heavy-tailed attention distribution), achieving higher compression with principled error control.

\subsection{Evaluation Metrics}

Most prior work evaluates KV cache compression using perplexity (PPL), which measures per-token prediction quality.
We argue that PPL systematically overestimates the quality impact of eviction-based compression, analogous to how PSNR overestimates quality degradation in video compression compared to perceptual metrics like VMAF~\citep{vmaf2016}.
This observation connects to a broader discussion in the compression community about task-aware distortion functions~\citep{blau2019rethinking}.

\subsection{Information-Theoretic Perspectives}

Shannon's rate-distortion theory~\citep{shannon1959} provides fundamental limits on lossy compression.
While widely applied in signal processing and video coding, its application to neural network intermediate representations (such as KV caches) is novel.
We show that the KV cache, after Hadamard rotation, is well-modeled as a Gaussian source, enabling direct application of the Gaussian rate-distortion function to derive compression limits.
The reverse water-filling algorithm~\citep{cover2006} for optimal bit allocation across sources with different variances provides the theoretical foundation for our four-tier token classification.
